\section{Introduction}

\subsection{The graph and the unexpected symmetry}

Let \(P\) denote the Petersen graph and let

\[
Y=L(P)
\]

be its line graph. The object studied in this paper is

\[
X=\KC(Y)\cong Y\times K_2,
\]

the canonical bipartite double cover of \(Y\).

The construction immediately displays two kinds of symmetry. Every
automorphism of \(Y\) lifts to \(X\), and the two sheets are exchanged by
the deck involution. Since

\[
\Aut(Y)\cong S_5,
\]

one naturally sees

\[
S_5\times C_2
\leq
\Aut(X).
\]

This visible subgroup is not the full automorphism group. The main result
is

\[
\Aut(X)\cong S_5\times S_3.
\]

The additional \(S_3\) is not an accidental enlargement discovered only
by enumeration. It is forced by a second realization of \(X\) as an
orbital graph on a homogeneous space of degree thirty.

\subsection{Main theorem}

The results of the paper may be summarized as follows.

\begin{theorem}
\label{thm:introduction-main}

Let

\[
H=V_{4,\mathrm{even}}
=
\{
1,
(12)(34),
(13)(24),
(14)(23)
\}
\leq S_5.
\]

Then the following statements hold.

\begin{enumerate}[label=\textup{(\roman*)}]

\item
The graph

\[
X=\KC(L(P))
\]

is connected, bipartite, and quartic, with

\[
|V(X)|=30,
\qquad
|E(X)|=60,
\qquad
\operatorname{girth}(X)=6.
\]

\item
Among the quartic self-paired orbital graphs of the action of \(S_5\) on

\[
S_5/H,
\]

exactly one is connected. That orbital graph is isomorphic to \(X\).

\item
The full automorphism group is

\[
\Aut(X)\cong S_5\times S_3.
\]

\item
The natural quotient frame has an automorphism orbit

\[
\frorbit
=
\Aut(X)\cdot\mathcal P_0
=
\{
\mathcal P_0,
\mathcal P_1,
\mathcal P_2
\}.
\]

Every member is a regular two-fold quotient frame with quotient
isomorphic to \(L(P)\). The three frames are pairwise block-disjoint and
contain forty-five distinct unordered vertex pairs.

\end{enumerate}
\end{theorem}

The theorem makes no claim that the displayed three-member orbit contains
every regular \(L(P)\)-quotient frame of \(X\).

\subsection{The source of the complementary \(S_3\)}

The second construction begins with the transitive action of \(S_5\) on

\[
\Omega=S_5/H.
\]

The centralizer of this action is determined by the standard coset-action
formula

\[
C_{\Sym(\Omega)}(S_5)
\cong
N_{S_5}(H)/H.
\]

Here

\[
N_{S_5}(H)\cong S_4,
\]

so the quotient is \(S_3\).

A centralizer element may permute the orbital relations rather than fix
them individually. In the present action, however, the certified quartic
orbital census contains exactly one connected member. Since relabeling
preserves connectedness and valency, that orbital relation is fixed by the
centralizer. The resulting \(S_3\) therefore acts by automorphisms of
\(X\).

This construction gives a subgroup

\[
S_5\times S_3
\leq
\Aut(X)
\]

of order \(720\). A certified full automorphism enumeration gives the
matching upper bound.

\subsection{Computational boundary}

Several finite statements in the paper are supplied by exact
certificates:

\begin{itemize}

\item the census of quartic self-paired orbitals on \(S_5/H\);

\item the graph isomorphism between the connected orbital and
\(\KC(L(P))\);

\item the full automorphism-group order;

\item the labeled alignment of the natural deck involution with the
complementary \(S_3\);

\item the explicit quotient-frame partitions and the forty-five-pair
census; and

\item the \(S_5\)-fixed binary cohomology calculation.

\end{itemize}

The proofs distinguish these certified inputs from conceptual deductions.
Section~9 records the evidence boundary and the reproduction contract.

\subsection{Organization}

Section~2 gives the required graph-theoretic, permutation-group, and
voltage-cover preliminaries. Section~3 constructs \(X\) and records its
elementary invariants. Section~4 develops the natural quotient frame.
Section~5 gives the independent homogeneous-space realization. Section~6
derives the full automorphism group from the centralizer action.
Section~7 interprets the complementary \(S_3\) through its quotient-frame
orbit. Section~8 places \(X\) in the invariant double-cover
classification. Section~9 records computational certification, and
Section~10 summarizes the result and the remaining open questions.
